\section{Evidence Map and Machine-Readable Table Reference}
% R005_APPENDIX_READER_FRAME
Reader frame. This appendix supports the main manuscript by explaining corpus role, evidence class, audit route, and limitation boundary. It is not a detached raw dump.


\label{app:oc133-source-to-pdf-trace}

This appendix is a reader-facing evidence map. It deliberately does not reproduce the raw provenance registers, source-path digests, or machine rows inside the monograph. Those complete registers remain in the version-pinned evidence corpus, where they can be checked by hash. The monograph records what each evidence class does for the scientific argument and how a reader should use it.


\subsection{Public Evidence Package Contents}

The public evidence manifest identifies the claim register, theorem registry, proof dependency graph, finite-model semantic checks, formalization inventory, retrospective replay prediction table, numeric replay QA table, domain evidence-boundary report, counterexample report, comparator and novelty registers, adversarial-review summary, and journal owner-review package index. The review package embeds the manifest and checksums; larger evidence files remain version-pinned public repository or corpus artifacts. The canonical filenames include: CLAIM\_LEDGER\_1\_3\_3.json, the corresponding public evidence-package record, FINITE\_MODEL\_CHECKS\_1\_3\_3.json, OC133V12.lean, the corresponding public evidence-package record, OC133\_TARGET\_BLIND\_PREDICTION\_TABLE.json, OC133\_NUMERIC\_REPLAY\_QA\_TABLE.json, the corresponding public evidence-package record, the corresponding public evidence-package record, OC\_1\_3\_3\_NOVELTY\_AND\_PRIORITY\_REGISTER.json, OC133\_LLM\_CERBERUS\_SUMMARY.json, SUBMISSION\_PACKAGE\_INDEX.json.


For exact SHA-256 values and public paths, use the corpus register, source-binding manifest, and checksums files for the same revision. Keeping those details in machine-readable artifacts avoids turning the scientific monograph into a raw hash catalogue while preserving reproducibility.


\subsection{Verification Summary}

\begin{itemize}[leftmargin=1.8em]

\item \textbf{Formalization inventory:} Lean-related files are cited as formalization inventory unless certificate binding and source-manifest checks are clean for the cited revision

\item \textbf{Finite semantic evidence:} 149 cases; any nonzero issue count is treated as a limitation and not as machine-checked closure

\item \textbf{Bounded replay lanes:} 5 lanes; promotion beyond replay support requires row-level comparator, negative-control, residual, and falsifier evidence

\item \textbf{Empirical promotion policy:} official snapshots are inputs, not empirical promotion by themselves

\item \textbf{Adversarial review:} open critical/high findings are recorded as 0/0 for the referenced review run

\item \textbf{Journal owner-review packets:} 8

\end{itemize}

\subsection{Publication-Grade Closure Rule}

A future release must fail if a primary scientific role is filled by an internal process memo, raw metadata wall, append-only surrogate, or verification package. The public manuscript must be a coherent text. Long registers may be cited and archived, but the public PDFs must explain how claims, proofs, data, review, and release boundaries fit together.
